% Proved: intrinsic recognition; four complete actual-file audits.
\subsection{Intrinsic small linear images of local permutation orbits}

\begin{lemma}[Proved: subgroup and point-stabilizer enumeration for small linear images]
\label{D:intrinsic-small-linear-permutation-images}
Fix a prime $p$ and a profinite group $\Lambda$ for which
\[
 |\operatorname{Hom}_{\rm cont}
       (\Lambda,\operatorname{GL}_D(\mathbb F_p))|
                      \le p^{aD^2}\qquad(D\ge1)          \tag{D.IM.1}
\]
with one constant $a$. For integers $N\ge2$ and $D\ge1$,
let $\mathcal T_{N,D}(\Lambda)$ be the isomorphism classes
of transitive $N$-element continuous $\Lambda$-sets whose
finite permutation image admits a faithful representation
over $\mathbb F_p$ of dimension at most $D$. Then
\[
 \log\max\{1,|\mathcal T_{N,D}(\Lambda)|\}
 \le \frac{2(\log p)^2}{\log2}D^4
                       +(a+2)(\log p)D^2.               \tag{D.IM.2}
\]
No linear representation or identification of the
permutation image is required as part of the given set.

More explicitly, there is a fixed collection of faithful
degree-$N$ permutation models, each embedded into
$\operatorname{GL}_D(\mathbb F_p)$, that contains every
action counted here. The logarithm of the maximum of
one and its number is at most
\[
            \frac{2(\log p)^2}{\log2}D^4
                              +2(\log p)D^2.             \tag{D.IM.3}
\]
Consequently, if positive integers $D_N$ satisfy
$D_N^4\le B N\log(N+1)$ for a fixed $B$, these
collections satisfy the hypotheses (D.MO.1)--(D.MO.2)
of Lemma~\ref{D:permutation-matrix-model-count}, after
enlarging its constant by an amount depending only
on $p$ and $B$.
\end{lemma}

\begin{proof}
Put $G=\operatorname{GL}_D(\mathbb F_p)$, $M=|G|$, and
$l=\lceil\log_2\max\{2,M\}\rceil$. Every subgroup
of a group of order at most $M$ has a generating tuple
of length at most $l$. To see this without assuming
any bound on its minimal generators, start with the
trivial subgroup and repeatedly adjoin an element
outside the subgroup already generated. Each proper
enlargement multiplies its order by at least two.
Padding the resulting tuple by identities gives
length exactly $l$.

Thus $G$ has at most $M^l$ subgroups. For each
$H\le G$, the same argument bounds the number of
subgroups $K\le H$ by $M^l$. Retain all ordered
pairs
\[
        (H,K),\qquad H\le G,\quad K\le H,\quad
        [H:K]=N,\quad \bigcap_{h\in H}hKh^{-1}=1.
                                                        \tag{D.IM.4}
\]
There are at most $M^{2l}$ such pairs. For each,
choose once a numbering of its left cosets $H/K$.
Left multiplication gives a faithful model
$H\hookrightarrow S_N$, and its other embedding
is the inclusion $H\hookrightarrow G$. These are
the promised fixed models.

Indeed, given a transitive set in the assertion,
write $H_0\le S_N$ for its actual permutation
image and $K_0$ for the stabilizer of a point.
The core of $K_0$ in $H_0$ is trivial because
$H_0$ acts faithfully. Choose any faithful
$\mathbb F_p$ representation of $H_0$ of
dimension at most $D$, and pad by a trivial
summand if necessary. Its images $H,K$ of
$H_0,K_0$ satisfy (D.IM.4). The original set
is isomorphic to the coset action of $H$ on
$H/K$, with the homomorphism $\Lambda\to H$
induced by that chosen representation.

In particular, the disjoint union over (D.IM.4)
of the actual sets
$\operatorname{Hom}_{\rm cont}(\Lambda,H)$
surjects onto $\mathcal T_{N,D}(\Lambda)$
after retaining its transitive actions. Its
cardinality is at most
\[
                  M^{2l}p^{aD^2},                       \tag{D.IM.5}
\]
by (D.IM.1) and the inclusion $H\le G$.
Choices of a faithful representation, point
stabilizer, coset numbering, or isomorphism
of the image are already covered by this
surjection; none needs an additional factor.
No injection of permutation isomorphism
classes into linear conjugacy classes is used.

If $M=1$, the collection (D.IM.4) is empty
for $N\ge2$, and the assertions follow.
If $M\ge2$, then $l\le\log M/\log2+1$
and $\log M\le D^2\log p$. Hence
\[
       2l\log M
       \le \frac{2}{\log2}(\log M)^2+2\log M
       \le \frac{2(\log p)^2}{\log2}D^4
                                      +2(\log p)D^2.    \tag{D.IM.6}
\]
This proves (D.IM.3), and (D.IM.5) proves
(D.IM.2). Finally $D_N^2\le D_N^4$ for
positive integers, so (D.IM.3) and
$D_N^2\log p\le B(\log p)N\log(N+1)$
give both matrix-model hypotheses.
The argument is elementary and uses no
classification of finite or primitive groups.
\end{proof}

\begin{corollary}[Proved: uncharged orbits with intrinsically small defining-characteristic images]
\label{D:intrinsic-linear-orbit-alternating-count}
In either Theorem~\ref{D:alternating-compressed-local-orbits}
or Theorem~\ref{D:original-boundary-matrix-alternating-count},
fix a sequence of positive integers $D_N$ with
$D_N^4\le B N\log(N+1)$. Enlarge the admissible
whole local orbits by allowing every degree-$N$
orbit whose finite image admits a faithful
$\mathbb F_p$ representation of dimension at
most $D_N$. Then all conclusions of the
respective theorem remain valid, with the
square-sum condition imposed only on the
remaining orbits.

In particular, for any fixed constants $c>0$
and $b\ge0$, all local orbit images of
faithful $\mathbb F_p$ dimension at most
$c(1+\log(N+1))^b$ may be left uncharged,
regardless of the particular transitive
action of that image. This is an intrinsic
condition on the whole local permutation
image, rather than a requirement that its
action be one prescribed affine or flag
action.
\end{corollary}

\begin{proof}
For complete puncture sources, (D.IM.1)
follows by multiplying the proved local
linear conjugacy-class count by the order
of the ambient linear group. For local
original-divisorial tame sources, it follows
from (D.OBM.4)--(D.OBM.5), which first
bound actual projective Hom sets and then
use $A\mapsto[\operatorname{diag}(A,1)]$.
For pro-prime-to-$p$ cyclic or trivial
sources the actual Hom bound follows
directly from one or zero generators.
Thus every one of the finitely many fixed
sources has an exponent $a$ as in
(D.IM.1).

Apply Lemma~\ref{D:intrinsic-small-linear-permutation-images}
at each degree, and adjoin its models to
the already chosen ones. The logarithm
of the number of models and the matrix
cost still satisfy (D.MO.2), with a new
fixed constant. The two cited theorems
therefore give the same factorial cover
bound, polynomial selected $A_n$ kernel
multiplicity, and finite generation of
the selected alternating simple product.
The whole local-group orbit remains the
unit of this condition.

For the final assertion take
$D_N=\max\{1,\lceil c(1+\log(N+1))^b\rceil\}$.
The ratio
\[
                  \frac{D_N^4}{N\log(N+1)}              \tag{D.IM.7}
\]
is bounded for $N\ge2$: its numerator
is bounded by a fixed power of
$1+\log(N+1)$, while its denominator
grows exponentially in $\log N$.
Hence a constant $B$ exists, depending
only on $c,b$, as required.
\end{proof}

\begin{remark}
This is a recognition enlargement within
the same matrix-model alternating criterion.
It does not assert that arbitrary local
permutation images have a small faithful
representation. No primitive-group
classification, hereditary cover estimate,
or unrestricted alternating finite-generation
claim is needed or inferred.
\end{remark}
